\begin{center}
    \textbf{ABSTRACT}
\end{center}

\vspace{0.25cm}

Colombian households often receive extensive and ambiguous school supply lists, forcing them to visit multiple stores to fulfill them—frequently without finding a single, comprehensive kit that meets their specific needs. Tukit processes these lists—whether provided as images, PDFs, or free-text—by comparing each item against a canonical catalog, utilizing both lexical and semantic similarity. It then assembles optimized kits using a multi-objective approach based on the NSGA-II algorithm, balancing factors such as cost, brand fidelity, product availability, and thematic preferences. The project followed a hybrid, encompassing methodology that integrates SCRUM, Design Thinking, RUP, and CRISP-DM, in alignment with IEEE standards 830, 1016, and 1058.
\begin{center}
    {\textbf{RESUMEN}}
\end{center}

\vspace{0.25cm}

Los hogares colombianos reciben listas escolares extensas y ambiguas, y deben recorrer
varios comercios para completarlas sin encontrar un kit que se ajuste a sus necesidades.
Tukit lee la lista en imagen, PDF o texto libre, compara cada renglón con un catálogo
canónico combinando el parecido léxico y el semántico. Posteriormente, arma kits con una optimización
multiobjetivo basada en NSGA-II que equilibra costo, fidelidad de marca, disponibilidad y
preferencias temáticas. El trabajo siguió una metodología híbrida contenedora que combina
SCRUM, Design Thinking, RUP y CRISP-DM con los estándares IEEE 830, 1016 y 1058.